\documentclass[UTF8,fontset=fandol,a4paper,openany,oneside]{ctexbook}
\usepackage{fontspec}
\setmonofont[Scale=0.80]{DejaVu Sans Mono}
\usepackage[margin=21mm,top=23mm,bottom=23mm]{geometry}
\usepackage{amsmath,amssymb,amsthm,mathtools}
\usepackage{booktabs,longtable,array}
\usepackage{xcolor,tikz}
\usetikzlibrary{arrows.meta,positioning}
\definecolor{InkBlue}{HTML}{19364D}
\definecolor{Teal}{HTML}{147D80}
\usepackage{xurl}
\usepackage[unicode,hidelinks]{hyperref}
\hypersetup{pdftitle={黎曼猜想：人类可读的证明主线},pdfauthor={基于指定Lean工程的中文数学重构},pdfsubject={ζ专用算术窗口、完整解析推导与深定理引用边界}}
\newtheorem{theorem}{定理}[chapter]
\newtheorem{lemma}[theorem]{引理}
\newtheorem{proposition}[theorem]{命题}
\newtheorem{inputtheorem}[theorem]{引用定理}
\theoremstyle{definition}
\newtheorem{definition}[theorem]{定义}
\theoremstyle{remark}
\newtheorem*{remark}{说明}
\renewcommand{\proofname}{证明}
\newcommand{\Rea}{\operatorname{Re}}
\newcommand{\Ima}{\operatorname{Im}}
\newcommand{\dd}{\,\mathrm d}
\newcommand{\Expl}[1]{\par\smallskip\noindent\textbf{这一步的意思。}#1\par\smallskip}
\setcounter{tocdepth}{1}
\setlength{\emergencystretch}{3em}
\setlength{\tabcolsep}{4pt}
\renewcommand{\arraystretch}{1.18}
\allowdisplaybreaks[1]
\raggedbottom
\urlstyle{tt}
\begin{document}
\begin{titlepage}
\vspace*{17mm}
\begin{center}
{\Large 从形式化证明到可以逐步检查的数学论证\par}
\vspace{16mm}
{\Huge\bfseries 黎曼猜想\par}
\vspace{8mm}
{\LARGE 人类可读的证明主线\par}
\vspace{12mm}
{\large 专门化到黎曼 $\zeta$ 函数\par}
{\large 算术窗口、极点与高阶导数、函数方程\par}
\end{center}
\vspace{16mm}
\noindent\rule{\textwidth}{0.5pt}
\par\medskip
本稿依据指定公开 Lean 工程及先前中文重构稿改写。正文完整展开从算术窗口估计到黎曼猜想的解析推导，并把决定性的算术几何输入准确列为引用定理；该输入的巨大上游证明没有在本稿中全部重写。
\par\medskip
\noindent\rule{\textwidth}{0.5pt}
\vfill
\noindent 整理日期：2026 年 11 月 22 日（北京时间）\par
\noindent 阅读基础：微积分、级数与复变函数的基本概念\par
\noindent 版本性质：第二版；增加完整源码汇编及算术接口展开；非作者原稿\par
\end{titlepage}
\frontmatter
\chapter*{怎样阅读这份证明}
\addcontentsline{toc}{chapter}{怎样阅读这份证明}
这份稿件把问题缩减为一个可以记住的矛盾：\textbf{如果临界线右边真的有零点，那么一个经过变换的函数就有极点；这个极点迫使某些高阶导数变大，而算术上构造出的好窗口却不允许它们那么大。}

正文不要求读者先理解整个一般 Dirichlet $L$ 函数理论。只取模 $1$ 的主特征，最终要研究的函数就是 $\zeta$。这会删去目标特征族的许多记号，但不会删去产生窗口的上游算术几何工作。

\textbf{本稿的证明范围必须准确理解：}第\ref{chap:arithmetic}章引用原工程的算术窗口定理；第\ref{chap:window}章完整证明该窗口怎样排除零点；第\ref{chap:finish}章完成参数选择与函数方程的推导。因此，本稿是一篇引用明确深定理的数学证明说明，\emph{不是}把全部上游引理都从基础重新证明的自足专著。把一项关键输入写成“引理”不会自动完成它的证明；它的来源、量词与本稿尚未展开的部分均在正文中列明。

截至本次查阅，Lean 与 Mathlib 社区于 2026 年 11 月 21 日发布的联合声明，已确认提交的命题与形式化证明，并报告命题比较与多个内核检查。该声明晚于前一版视频所采用的 MIT 12:55 UTC 状态快照。此项外部核验与本稿的人类文字重构是两件分别记录的工作；本次没有重新运行全工程核验。源码汇编恢复的是原项目各模块的文本；模块哈希一致性核验与 Lean 内核核验分别记录。原作者的四页公告也说明，论证很长，仍需进一步整理、理解和简化。参见文末资料[1]--[3]。

阅读时，可以先读每章的解释段和最后一章，掌握全链条；再回到第\ref{chap:window}章检查真正承担矛盾的每个不等式。全文不把有限个数值零点当作证明，也不把单侧上界偷换成绝对值上界。
\tableofcontents
\mainmatter
